Un radiotélescope dont l'antenne parabolique a un diamètre de \qty{\v{Dm}}{m} observe le rayonnement de longueur d'onde \qty{\v{lamr}}{cm} émis par la galaxie d'Andromède, située à \qty{2.4e22}{m}.

\begin{enumerate}[label=\alph*)]
    \item Déterminez l'angle limite de résolution du radiotélescope, en degrés.
    \item Quelle est la taille du plus petit détail qu'il peut distinguer dans la galaxie d'Andromède ?
    \item Quel diamètre aurait un télescope optique qui a la même résolution pour une lumière de \qty{550}{nm} ? Expliquez pourquoi les radiotélescopes sont si grands.
\end{enumerate}
